\section{El Patrón MVC y su Implementación en Frameworks Modernos}
\subsection{Historia y evolución del MVC}

El patrón Modelo-Vista-Controlador (MVC) fue descrito por primera vez por \textbf{Trygve Reenskaug en 1979}
mientras trabajaba en el entorno \textbf{Smalltalk-80} del Xerox PARC. En su formulación original, el
\textit{modelo} representa el estado y la lógica de negocio, la \textit{vista} es la representación visual de ese
estado y el \textit{controlador} es el componente que maneja la entrada del usuario (eventos de teclado y ratón).
El objetivo era dividir responsabilidades para permitir que múltiples vistas observaran un mismo modelo y que los
cambios en los datos se reflejaran automáticamente en la interfaz.

Cuando el MVC se trasladó al entorno web, la semántica de los componentes cambió para adaptarse a la
arquitectura \textit{request/response} del protocolo HTTP. En el modelo web, el controlador ya no maneja eventos de
entrada locales, sino que procesa \textbf{solicitudes HTTP}: recibe un verbo y una URI, invoca al modelo y
selecciona una vista. La vista puede renderizarse en el servidor mediante un motor de plantillas o en el cliente
mediante una aplicación de una sola página (SPA). Por su parte, el modelo se amplía para incluir repositorios,
servicios de aplicación y objetos de acceso a datos (DAO). De esta forma, MVC se convierte en el patrón
arquitectónico dominante de la mayoría de frameworks modernos.

Sobre el MVC clásico se desarrollaron variantes que resuelven limitaciones concretas. El patrón \textbf{MVP}
(Model-View-Presenter) elimina la referencia directa entre la vista y el modelo: el presentador actúa como
intermediario único y la vista es pasiva. El patrón \textbf{MVVM} (Model-View-ViewModel) introduce un
\textit{ViewModel} con enlace bidireccional de datos (\textit{data binding}), muy utilizado en aplicaciones SPA
como Angular, donde una directiva como \texttt{[(ngModel)]} sincroniza automáticamente la vista y el modelo.

\subsection{El patrón Front Controller}

Todos los frameworks MVC web relevantes (Laravel, Symfony, Spring MVC, ASP.NET Core, Django) emplean el patrón
\textbf{Front Controller}: un punto de entrada único que centraliza el manejo de errores, la autenticación, el
logging y el enrutamiento. En Spring Boot, ese punto de entrada es \texttt{DispatcherServlet}, declarado y
configurado automáticamente por el arranque de la aplicación. Este diseño evita repetir lógica transversal
(auth, CORS, cabeceras de seguridad) en cada controlador.

El ciclo completo de una petición MVC en Spring Boot, aplicado al PFC, es el siguiente:

\begin{enumerate}[label=\textbf{Paso \arabic*.}]
    \item \textbf{Recepción de la petición HTTP:} el cliente envía, por ejemplo, \texttt{GET /api/conductores}
    con el token JWT. El contenedor embebido (Apache Tomcat) recibe la solicitud.
    \item \textbf{Filtros de seguridad:} \texttt{JwtAuthenticationFilter} se ejecuta antes de
    \texttt{UsernamePasswordAuthenticationFilter} (ver \texttt{SecurityConfig.java}) y, si el token es válido y no
    está en la \textit{blacklist} de Redis, autentica la petición sin estado.
    \item \textbf{Enrutamiento:} \texttt{DispatcherServlet} consulta el \texttt{HandlerMapping} y localiza el
    \texttt{ConductorController.listar()} que atiende \texttt{GET /api/conductores}.
    \item \textbf{Invocación del controlador:} el \texttt{HandlerAdapter} invoca el método del controlador,
    que recibe un \texttt{Pageable} resuelto automáticamente por Spring Data.
    \item \textbf{Interacción con el modelo:} el controlador delega en \texttt{ConductorService.listar()}, que a su
    vez consulta \texttt{ConductorRepository} (capaz de usar la caché de Redis vía \texttt{@Cacheable}) y devuelve
    un \texttt{Page<ConductorResponse>}.
    \item \textbf{Serialización de la respuesta:} el \texttt{HttpMessageConverter} (Jackson) serializa el objeto de
    respuesta a JSON y el servidor devuelve el código HTTP 200.
\end{enumerate}

Este flujo se evidencia en el siguiente fragmento real del controlador:

\begin{lstlisting}[style=javaStyle, label=lst:c1, caption={ConductorController.java — endpoint GET /api/conductores}]
@RestController
@RequestMapping("/api/conductores")
@RequiredArgsConstructor
public class ConductorController {

    private final ConductorService conductorService;

    @GetMapping
    public ResponseEntity<Page<ConductorResponse>> listar(
            @PageableDefault(size = 10, sort = "id") Pageable pageable
    ) {
        return ResponseEntity.ok(conductorService.listar(pageable));
    }
}
\end{lstlisting}

\subsection{Comparativa de frameworks MVC en 2025}

Para justificar la elección del framework del PFC se compararon los cinco frameworks MVC de mayor adopción en el
mercado latinoamericano (Tabla~\ref{tab:frameworks}). La comparación considera el paradigma, el ORM incluido, el
motor de plantillas, el ecosistema de paquetes y el rendimiento reportado en \textit{TechEmpower Framework
Benchmarks} (Round 22).

\begin{table}[H]
\centering
\small
\setlength{\tabcolsep}{5pt}
\begin{tabular}{@{}p{2.6cm}p{3.4cm}p{2.6cm}p{2.6cm}p{3.2cm}@{}}
\toprule
\textbf{Framework} & \textbf{Paradigma} & \textbf{ORM incluido} & \textbf{Motor de plantillas} & \textbf{Ecosistema / mercado} \\
\midrule
Laravel 11.x & Convención sobre configuración & Eloquent ORM & Blade & Amplio en LATAM, Packagist \\
Symfony 7.x & Explícito y extensible & Doctrine ORM & Twig & Gran adopción empresarial \\
Spring Boot 3.x & Configuración por anotaciones, opinionated & Spring Data JPA & Thymeleaf / SPA & Muy amplio, Java dominante \\
ASP.NET Core 8.x & Convención + C\# fuertemente tipado & EF Core & Razor & Medio en LATAM \\
Django 5.x & Convención sobre configuración & Django ORM & DTL / Jinja2 & Medio, buena en Sudamérica \\
\bottomrule
\end{tabular}
\caption{Frameworks MVC evaluados para el PFC (adaptado de documentación oficial de cada framework).}
\label{tab:frameworks}
\end{table}

En el benchmark TechEmpower Round 22, los frameworks JVM (Spring Boot, Vert.x, Quarkus) se ubican de forma
consistente en la franja media-alta de rendimiento, con latencias muy competitivas cuando se combinan con caché.
Para el PFC se eligió \textbf{Spring Boot} por su ecosistema maduro de seguridad (\texttt{spring-security}),
soporte de primera clase a JPA/PostgreSQL, integración natural con Redis y por el conocimiento previo del equipo
documentado en el \textbf{ADR-001} (arquitectura monolítica vs. microservicios). La decisión técnica queda
registrada en \texttt{docs/adr/adr-001-arquitectura-monolitica.md}.